\chapter[CAPÍTULO VII]{CONCLUSIONES}

\section{Conclusiones respecto a los objetivos del Proyecto de Mejora / Innovación}

\subsection{Resultados}

La implementación del Sistema Web Help-Desk Centralizado en J\&P Periféricos S.A.C. permitió alcanzar los objetivos planteados para la gestión del servicio de soporte técnico y el control de inventario de activos TI orientados a empresas clientes externas. Los resultados operativos y de validación funcional y técnica alcanzados se resumen a continuación, alineados con cada objetivo específico:

\begin{itemize}[leftmargin=1.27cm]
    \item \textbf{Recepción y asignación centralizada de solicitudes:} se eliminó la dependencia de canales informales (llamadas telefónicas y WhatsApp sin registro) mediante un módulo de gestión de tickets con código PIN único y flujo de estados auditable, garantizando la trazabilidad del ciclo de vida de cada requerimiento desde su apertura hasta su cierre, asignación y escalamiento.

    \item \textbf{Expediente e inventario multinivel de activos TI:} se implementó un módulo de gestión de expedientes e inventarios de TI (hardware y software) parametrizado por empresa cliente, sucursal, área y trabajador, que provee de manera inmediata la hoja de vida histórica, los componentes instalados y las licencias de software por equipo, reduciendo la búsqueda manual en hojas de cálculo dispersas y desactualizadas.

    \item \textbf{Coordinación de visitas \textit{in-situ} sin cruces horarios:} se integró un módulo de agendamiento y programación de citas presenciales con motor de autoasignación por zonas geográficas y control anti-cruces de franjas horarias, agrupando las atenciones presenciales de manera eficiente y reduciendo tiempos muertos, traslados duplicados y reprocesos en la reprogramación.

    \item \textbf{Control digital de planes, contratos y SLA:} se desarrolló un módulo que valida de forma automatizada la vigencia de los planes contratados, la cantidad máxima de equipos cubiertos y los tiempos máximos de atención comprometidos, asegurando el cumplimiento formal de los Acuerdos de Nivel de Servicio (SLA) y evitando la atención de activos fuera de cobertura.

    \item \textbf{Flujo estandarizado de derivación y escalamiento:} se estandarizó el procedimiento de derivación y escalamiento de tickets, con registro del motivo de reasignación y coordinación con casos recurrentes, asegurando el seguimiento continuo y la resolución oportuna de las atenciones a nivel técnico.

    \item \textbf{Tableros de métricas y KPIs:} se construyó un módulo de tableros de control con indicadores clave de rendimiento (MTTA, MTTR, FCR y CSAT) y dimensiones de calidad de servicio, permitiendo a la gerencia evaluar objetivamente los tiempos de respuesta y resolución, el cumplimiento de SLA y la satisfacción del cliente.
\end{itemize}

En términos técnicos, el sistema se desarrolló con un backend en NestJS y Prisma ORM sobre MySQL para los datos transaccionales y relacionales, MongoDB para el historial de chats e interacciones, Redis con \textit{Pub/Sub} para eventos en tiempo real y WebSockets para la comunicación bidireccional; el frontend se implementó en Vite + React con \textit{shadcn/ui} y arquitectura modular por roles (ADMINISTRADOR, SOPORTE\_TECNICO, SOPORTE\_INSITU, CLIENTE\_EMPRESA, CLIENTE\_SUCURSAL y CLIENTE\_TRABAJADOR). La arquitectura multiinquilino (multi-tenant) garantiza el aislamiento de datos por empresa cliente, mientras que la jerarquía multinivel (empresa, sucursal, área, trabajador) permite la trazabilidad precisa de cada activo y de cada incidencia.

\subsection{Discusión}

Los resultados obtenidos coinciden con la tendencia reportada por los antecedentes nacionales e internacionales que evidencian que la adopción de sistemas web estructurados para la gestión de incidencias reduce los tiempos de atención y mejora la trazabilidad operativa. Mendoza Lizana (2023) reportó una relación positiva muy alta del 93,7\,\% entre ITIL v.4 y la gestión de incidencias; Figueroa Santander (2022) evidenció un incremento del promedio de incidencias resueltas del 56,52\,\% al 70,13\,\% y una reducción del tiempo promedio de resolución de 38,36 a 12,58 minutos; y Villalva Mendivil (2022) demostró que un sistema \textit{Helpdesk} estructurado incrementó la atención de incidencias del 63,74\,\% al 88,68\,\% y la resolución en primer nivel del 48,52\,\% al 93,57\,\%. De manera similar, Táez Valdiviezo (2021) reportó una mejora promedio del 25,08\,\% en los tiempos de atención al contrastar registros manuales frente a una herramienta \textit{cloud Help Desk} con métricas ITIL 4, y Lanuza Bustamante (2024) evidenció, mediante la prueba U de Mann-Whitney ($\text{sig.} = 0,000 < 0,05$), una reducción del 77,76\,\% en el tiempo de incidencias y del 73,9\,\% en la cantidad de incidencias tras implementar prácticas basadas en ITIL 4.0.

En el ámbito de las plataformas web, Velasco Castro (2023) reportó una usabilidad del 95\,\% bajo la norma ISO/IEC 9126 en un sistema web Help-Desk que sustituyó registros manuales, y Teixeira Ferreira y Knoll Ricardo (2026) evidenciaron una reducción de hasta el 40\,\% en el tiempo de clasificación de solicitudes al centralizar la atención en una plataforma web propia. Díaz Hidalgo (2024) demostró la viabilidad de migrar la mesa de ayuda a la nube conservando el historial de tickets y disponiendo de tableros estadísticos de monitoreo, mientras que Gonzales Delgado (2025) registró que, tras implementar un sistema Help-Desk basado en ITIL 4, las incidencias resueltas aumentaron a 432 (+18\,\%), las no resueltas cayeron a 32 (-67\,\%) y el tiempo de resolución bajó a 178 horas (-33\,\%), con calidad alta del 52,2\,\% y eliminación del nivel bajo (0,0\,\%).

Respecto a la trazabilidad y el control de activos, Chávez Cotera (2022) destacó que la trazabilidad del parque informático es un requisito para mejorar la gestión de incidencias, al evitar diagnósticos repetitivos y reescritura de informes, y Vargas Calapucha y Romero Licuy (2026) evidenciaron que la centralización del inventario de equipos tecnológicos en una plataforma web optimiza la organización del mantenimiento, la seguridad de los datos y la eficiencia operativa. Recalde Manrique (2022) respaldó la gestión de activos de TI bajo ITIL 4 (ciclo de vida, licencias y reportes) mediante un prototipo web con Vue.js, Spring y MySQL, validado al 100\,\% por expertos y usuarios, aunque sin métricas cuantitativas de mejora (antes/después).

Las diferencias de alcance respecto a los antecedentes son las siguientes:

\begin{itemize}[leftmargin=1.27cm]
    \item \textbf{Soporte interno vs. atención a empresas clientes externas multinivel:} la mayoría de los antecedentes (Fonseca Calcin, 2025; Villalva Mendivil, 2022; Chávez Cotera, 2022; Figueroa Santander, 2022; Mendoza Lizana, 2023; Lanuza Bustamante, 2024; Torres Hernández, 2023) se centran en el soporte técnico interno de una organización (universidades, empresas privadas o entidades públicas). El presente proyecto, en cambio, atiende a empresas clientes externas con múltiples sedes y trabajadores, incorporando una jerarquía multinivel (empresa, sucursal, área, trabajador) y roles diferenciados (CLIENTE\_EMPRESA, CLIENTE\_SUCURSAL, CLIENTE\_TRABAJADOR), lo que amplía el alcance funcional y la trazabilidad del servicio.

    \item \textbf{Inventario comercial vs. ITAM de soporte:} Chávez Palomino (2024) validó cuantitativamente (23,48\,\% en control de inventario, 21,65\,\% en exactitud, 25,32\,\% en rotación) la mejora del inventario mediante un sistema web, pero en un contexto de inventario comercial de productos de seguridad industrial (rotación, exactitud de stock). El presente proyecto aplica el control de inventario a activos TI (hardware/software, licencias, mantenimientos) de empresas clientes externas, con expediente técnico por activo y trazabilidad de cambios, lo que corresponde a una gestión ITAM de soporte y no a un inventario de venta.

    \item \textbf{Validación de prototipo vs. operación:} Recalde Manrique (2022) y Torres Hernández (2023) validaron prototipos o sistemas en contextos acotados (validación por expertos y usuarios, o atención a 60 registros en tres meses en una empresa del rubro pesca). El presente proyecto integra en una sola plataforma la gestión de tickets, el inventario multinivel, el agendamiento de visitas \textit{in-situ}, la validación de SLA y los tableros de KPIs, orientada a la operación continua con empresas clientes externas.
\end{itemize}

El aporte principal del proyecto reside en unificar en una sola plataforma web centralizada la gestión del servicio de soporte técnico y el control de inventario de activos TI para clientes corporativos externos, resolviendo de forma integrada las causas raíz CA1 a CA7 identificadas en el diagnóstico (recepción informal sin trazabilidad, inventario disperso en Excel, descoordinación de visitas \textit{in-situ}, escalamiento sin flujo estandarizado, verificación manual de SLA, reportes manuales y carencia de KPIs), y aportando trazabilidad auditable, métricas objetivas y cumplimiento automatizado de los acuerdos de servicio.

\subsection{Conclusión}

Se concluye que el desarrollo e implementación del Sistema Web Help-Desk Centralizado para la gestión del servicio de soporte técnico y el control de inventario de activos TI en J\&P Periféricos S.A.C. responde de manera integral al problema principal identificado: la desorganización y falta de trazabilidad centralizada en la gestión de solicitudes de soporte técnico y registro de inventario informático multinivel. La solución reemplaza los canales informales de recepción (llamadas y WhatsApp sin registro) por un flujo de tickets estandarizado con código PIN único y estados auditables; centraliza el expediente e inventario de activos TI (hardware y software) por empresa, sucursal, área y trabajador; coordina las visitas \textit{in-situ} sin cruces de agenda mediante autoasignación por zona geográfica; valida de forma automatizada los planes, contratos y SLA de cada cliente corporativo; estandariza la derivación y el escalamiento de tickets; y provee tableros de control con KPIs (MTTA, MTTR, FCR y CSAT) para la evaluación objetiva del desempeño técnico.

En consecuencia, los objetivos específicos del proyecto se cumplieron: se automatizó la recepción y asignación de solicitudes de soporte; se implementó el módulo de expedientes e inventarios de TI multinivel; se integró el agendamiento de citas presenciales; se desarrolló el control digital de planes y SLA; se estandarizó el flujo de derivación y escalamiento; y se construyeron los tableros de métricas para la evaluación del servicio. De este modo, la solución aporta trazabilidad, eficiencia operativa y una base objetiva para la toma de decisiones gerenciales en el servicio de soporte técnico de J\&P Periféricos S.A.C., orientado a sus empresas clientes externas y sus activos TI.